\chapter{Modelagem, Relat\'orios e Resultados}
\label{cap:resultados}

\section{Modelos candidatos e crit\'erios de sele\c{c}\~ao}
\label{sec:candidatos}

A etapa de modelagem comparou candidatos de regress\~ao, candidatos de classifica\c{c}\~ao e baselines determin\'isticos. Essa organiza\c{c}\~ao foi importante porque o objetivo n\~ao era apenas treinar um algoritmo sofisticado, mas verificar se ele realmente acrescentava valor em rela\c{c}\~ao a regras simples baseadas no pr\'oprio hist\'orico de notas.

Os candidatos avaliados foram:

\begin{itemize}
    \item \textit{RandomForestRegressor} e \textit{RandomForestClassifier};
    \item \textit{HistGradientBoostingRegressor} e \textit{HistGradientBoostingClassifier};
    \item baseline de \'ultima nota;
    \item baseline de m\'edia das duas \'ultimas notas;
    \item baseline de m\'edia das tr\^es \'ultimas notas;
    \item baseline h\'ibrido.
\end{itemize}

Para regress\~ao, o crit\'erio principal de sele\c{c}\~ao foi o \sigla{MAE}{Mean Absolute Error}. Para classifica\c{c}\~ao de risco, o crit\'erio principal foi o \sigla{F1}{F1-score}, acompanhado de \textit{precision} e \textit{recall}. O MAE foi escolhido porque expressa o erro diretamente em pontos de nota, o que facilita a interpreta\c{c}\~ao pedag\'ogica. O F1 foi usado porque equilibra a capacidade de encontrar casos cr\'iticos com a precis\~ao dos alertas emitidos.

\section{Compara\c{c}\~ao dos candidatos na valida\c{c}\~ao}
\label{sec:comparacao-candidatos}

A Tabela~\ref{tab:validacao-regressao} mostra a compara\c{c}\~ao dos principais candidatos de regress\~ao na etapa de valida\c{c}\~ao. No modo de previs\~ao de nota, o \textit{Random Forest} foi ligeiramente superior ao \textit{HistGradientBoosting} em MAE, por isso foi selecionado como regressor principal desse modo. No modo de alerta de risco, por outro lado, o \textit{HistGradientBoosting} foi o melhor regressor dentro do recorte com hist\'orico m\'inimo maior.

\begin{table}[ht]
\caption{Compara\c{c}\~ao dos candidatos de regress\~ao na valida\c{c}\~ao.}
\label{tab:validacao-regressao}
\centering
\footnotesize
\begin{tabular}{l c c c}
\hline
\textbf{Modo e candidato} & \textbf{MAE} & \textbf{RMSE} & \textbf{Erro $\leq$ 0,5} \\
\hline
Previs\~ao -- Random Forest & 0,8152 & 1,0907 & 0,4261 \\
Previs\~ao -- HistGradientBoosting & 0,8191 & 1,0926 & 0,4156 \\
Previs\~ao -- baseline de \'ultima nota & 0,9148 & 1,3015 & 0,4531 \\
Alerta -- HistGradientBoosting & 0,8124 & 1,1113 & 0,4429 \\
Alerta -- Random Forest & 0,8327 & 1,1167 & 0,4174 \\
Alerta -- baseline de duas notas & 0,9196 & 1,2746 & 0,4248 \\
\hline
\end{tabular}
\fontetab{Elaborada pelo autor com base nos artefatos da pipeline.}
\end{table}

A Tabela~\ref{tab:validacao-classificacao} apresenta a compara\c{c}\~ao dos candidatos de classifica\c{c}\~ao para risco. O resultado mostra duas evid\^encias importantes. Primeiro, baselines simples ainda s\~ao competitivos e precisam ser considerados. Segundo, no modo especificamente calibrado para alerta de risco, o \textit{HistGradientBoostingClassifier} superou os demais candidatos em F1.

\begin{table}[ht]
\caption{Compara\c{c}\~ao dos candidatos de classifica\c{c}\~ao na valida\c{c}\~ao.}
\label{tab:validacao-classificacao}
\centering
\footnotesize
\begin{tabular}{l c c c}
\hline
\textbf{Modo e candidato} & \textbf{Precision} & \textbf{Recall} & \textbf{F1} \\
\hline
Previs\~ao -- baseline de duas notas & 0,7011 & 0,7744 & 0,7359 \\
Previs\~ao -- HistGradientBoosting & 0,6172 & 0,8937 & 0,7302 \\
Previs\~ao -- Random Forest & 0,5743 & 0,9158 & 0,7059 \\
Alerta -- HistGradientBoosting & 0,6980 & 0,8255 & 0,7564 \\
Alerta -- baseline de duas notas & 0,7165 & 0,7908 & 0,7518 \\
Alerta -- Random Forest & 0,6091 & 0,9036 & 0,7277 \\
\hline
\end{tabular}
\fontetab{Elaborada pelo autor com base nos artefatos da pipeline.}
\end{table}

\section{Resultados do modo de previs\~ao de nota}
\label{sec:prev-nota}

No modo \texttt{previsao\_nota}, o melhor cen\'ario foi obtido com hist\'orico m\'inimo igual a 1. A base de modelagem alcan\c{c}ou 106.721 linhas, 1.378 alunos \'unicos, 53 disciplinas e 6 anos letivos considerados. O conjunto temporal foi organizado com treino em anos anteriores, valida\c{c}\~ao em 2024 e teste em 2025.

O melhor candidato de regress\~ao foi o \textit{RandomForestRegressor}, com os resultados finais apresentados na Tabela~\ref{tab:prevnota}. Em termos pr\'aticos, isso significa que a previs\~ao da pr\'oxima nota ficou, em m\'edia, a aproximadamente 0,77 ponto da nota realmente observada, superando a heur\'istica simples de repetir a nota anterior.

\begin{table}[ht]
\caption{Resultados finais do modo de previs\~ao de nota no teste temporal.}
\label{tab:prevnota}
\centering
\begin{tabular}{l c}
\hline
\textbf{M\'etrica} & \textbf{Valor} \\
\hline
MAE do modelo & 0,7659 \\
RMSE do modelo & 1,0519 \\
Acerto com erro $\leq$ 0,5 & 0,4649 \\
MAE do baseline de \'ultima nota & 0,9138 \\
RMSE do baseline de \'ultima nota & 1,2967 \\
\hline
\end{tabular}
\fontetab{Elaborada pelo autor com base nos artefatos da pipeline.}
\end{table}

\section{Resultados do modo de alerta de risco}
\label{sec:alerta-risco}

No modo \texttt{alerta\_risco}, o melhor cen\'ario foi obtido com hist\'orico m\'inimo igual a 2. Nessa configura\c{c}\~ao, a base de modelagem alcan\c{c}ou 50.065 linhas, 1.179 alunos \'unicos e 53 disciplinas. O melhor classificador foi o \textit{HistGradientBoostingClassifier}, com os resultados apresentados na Tabela~\ref{tab:riscos}.

\begin{table}[ht]
\caption{Resultados finais do modo de alerta de risco no teste temporal.}
\label{tab:riscos}
\centering
\begin{tabular}{l c}
\hline
\textbf{M\'etrica} & \textbf{Valor} \\
\hline
Precision do modelo & 0,6442 \\
Recall do modelo & 0,8794 \\
F1 do modelo & 0,7436 \\
F1 do baseline & 0,7243 \\
\hline
\end{tabular}
\fontetab{Elaborada pelo autor com base nos artefatos da pipeline.}
\end{table}

O destaque est\'a no \textit{recall}. Em um sistema de alerta, perder casos realmente cr\'iticos tende a ser mais grave do que sinalizar alguns casos adicionais para acompanhamento. Por isso, esse resultado foi considerado adequado ao objetivo pedag\'ogico do trabalho: priorizar alunos que merecem revis\~ao humana, n\~ao produzir uma decis\~ao autom\'atica.

A Tabela~\ref{tab:matriz-risco} resume a matriz de confus\~ao do modo de alerta de risco no teste temporal. Dos 2.794 casos que realmente tiveram pr\'oxima nota abaixo de 6,0, o modelo recuperou 2.457. Por outro lado, tamb\'em sinalizou 1.357 casos que n\~ao se confirmaram como risco no alvo observado. Esse comportamento \'e coerente com uma ferramenta de triagem, pois privilegia cobertura dos casos cr\'iticos.

\begin{table}[ht]
\caption{Matriz de confus\~ao do alerta de risco no teste temporal.}
\label{tab:matriz-risco}
\centering
\begin{tabular}{l c c}
\hline
\textbf{Risco real} & \textbf{Predito sem risco} & \textbf{Predito com risco} \\
\hline
Sem risco & 9.092 & 1.357 \\
Com risco & 337 & 2.457 \\
\hline
\end{tabular}
\fontetab{Elaborada pelo autor com base nos artefatos da pipeline.}
\end{table}

\section{Comparando hist\'oricos m\'inimos}
\label{sec:minhistory}

Os testes de hist\'orico m\'inimo mostraram um resultado relevante para a arquitetura do projeto:

\begin{itemize}
    \item \texttt{min\_history = 1} foi melhor para regress\~ao da pr\'oxima nota, pois preservou mais alunos e disciplinas na base;
    \item \texttt{min\_history = 2} foi melhor para classifica\c{c}\~ao de risco, pois exigiu maior evid\^encia hist\'orica antes de sinalizar alerta;
    \item \texttt{min\_history = 3} reduziu demais a quantidade de casos e inviabilizou a valida\c{c}\~ao temporal na rodada analisada.
\end{itemize}

Essa evid\^encia levou \`a separa\c{c}\~ao da pipeline em dois modos complementares de uso, cada um com sua configura\c{c}\~ao mais adequada.

\section{Exemplo interpretado de alerta}
\label{sec:exemplo-alerta}

A Tabela~\ref{tab:exemplo-alerta} apresenta um caso pseudonimizado retirado dos artefatos de teste. O exemplo \'e \'util porque mostra por que o projeto manteve um classificador de risco, e n\~ao apenas uma regra baseada na nota prevista.

\begin{table}[ht]
\caption{Exemplo pseudonimizado de caso sinalizado pelo alerta de risco.}
\label{tab:exemplo-alerta}
\centering
\footnotesize
\begin{tabular}{p{4.9cm} p{8.5cm}}
\hline
\textbf{Campo} & \textbf{Valor observado} \\
\hline
Aluno, s\'erie e disciplina & Aluno 01790, 1\textordfeminine{} S\'erie, Biologia 2 \\
Etapa analisada & 3\textordmasculine{} bimestre de 2025 \\
Nota atual e pr\'oxima nota real & 6,3 na etapa atual; 5,8 na etapa seguinte \\
Notas anteriores na disciplina & 5,5 e 4,5 \\
M\'edia das duas \'ultimas notas & 5,0 \\
Faltas acumuladas & 7 faltas at\'e a etapa analisada \\
Pagamentos pendentes no ano & 0 pend\^encias registradas \\
M\'edia da coorte na etapa & 6,98 \\
Professor vinculado & Professor 00019 \\
Nota prevista pelo regressor & 6,58 \\
Alerta do classificador & Risco sinalizado corretamente \\
\hline
\end{tabular}
\fontetab{Elaborada pelo autor com base nos artefatos da pipeline.}
\end{table}

Se o sistema usasse apenas a regra ``nota prevista menor que 6'', esse caso n\~ao seria sinalizado, pois a previs\~ao num\'erica foi 6,58. O classificador, entretanto, identificou o risco a partir do conjunto de sinais: hist\'orico anterior fraco na disciplina, m\'edia recente abaixo da linha de seguran\c{c}a, faltas acumuladas e desempenho abaixo da m\'edia da coorte. Esse exemplo mostra que a regress\~ao e a classifica\c{c}\~ao s\~ao complementares.

\section{Gera\c{c}\~ao de relat\'orios escolares}
\label{sec:relatorios}

Os resultados preditivos n\~ao foram o ponto final da solu\c{c}\~ao. A partir deles, o projeto passou a gerar relat\'orios para tr\^es perfis institucionais:

\begin{itemize}
    \item professor: com nota atual, nota prevista, fatores de alerta e recomenda\c{c}\~ao pedag\'ogica;
    \item coordena\c{c}\~ao: com vis\~ao por s\'erie, disciplina, ranking executivo e consolidados por professor;
    \item secretaria: com sinais administrativos relacionados a frequ\^encia, financeiro e prioridade de contato.
\end{itemize}

Tamb\'em foram constru\'idos relat\'orios por professor, incluindo consolidados por turma e consolidados gerais, com contagem de alunos em risco baixo, moderado e alto. Essa camada aproxima a modelagem do cotidiano escolar, pois transforma erro, probabilidade e predi\c{c}\~ao em informa\c{c}\~ao compreens\'ivel por cada perfil de uso.

\section{An\'alise de erro e calibra\c{c}\~ao}
\label{sec:erro}

A an\'alise de erro mostrou maior dificuldade de previs\~ao em disciplinas como Matem\'atica 3, Matem\'atica 4 e 5, Biologia 2, Hist\'oria 2 e Qu\'imica 3. Tamb\'em houve maior erro nas faixas de nota abaixo de 7, especialmente porque esses casos s\~ao mais inst\'aveis e pedagogicamente mais sens\'iveis.

\begin{table}[ht]
\caption{Principais pontos de erro no modo de alerta de risco.}
\label{tab:erro-alerta}
\centering
\footnotesize
\begin{tabular}{l c c c}
\hline
\textbf{Grupo analisado} & \textbf{Amostras} & \textbf{MAE} & \textbf{Acur\'acia de risco} \\
\hline
Matem\'atica 4 e 5 & 81 & 1,9489 & 0,7778 \\
Biologia 2 & 260 & 1,5895 & 0,6500 \\
Hist\'oria 2 & 81 & 1,3741 & 0,6667 \\
Faixa abaixo de 6 & 2.794 & 1,0395 & 0,8794 \\
Faixa de 6 a 6,99 & 1.405 & 0,9718 & 0,5103 \\
\hline
\end{tabular}
\fontetab{Elaborada pelo autor com base nos artefatos da pipeline.}
\end{table}

Na camada operacional, a primeira regra de risco gerava excesso de casos classificados como alto risco. Por isso, foi criada uma pontua\c{c}\~ao composta que passou a combinar:

\begin{itemize}
    \item predi\c{c}\~ao de risco do modelo;
    \item nota prevista;
    \item nota atual;
    \item varia\c{c}\~ao esperada;
    \item faltas acumuladas;
    \item pend\^encias financeiras.
\end{itemize}

Com essa calibra\c{c}\~ao, a prioriza\c{c}\~ao ficou mais compat\'ivel com o uso institucional, reduzindo sobrecarga sobre coordena\c{c}\~ao e professores. Assim, a sa\'ida final do sistema n\~ao deve ser lida como senten\c{c}a sobre o aluno, mas como convite estruturado \`a investiga\c{c}\~ao pedag\'ogica.
